\documentclass[12pt,a4paper]{article}
\usepackage[top=22mm,bottom=24mm,left=27mm,right=27mm]{geometry}
\usepackage{amsmath,amssymb}
\usepackage{xcolor}
\definecolor{link}{HTML}{205973}
\usepackage[hypertexnames=false,colorlinks=true,linkcolor=link,urlcolor=link]{hyperref}
\hypersetup{pdftitle={Theorem 1. An approximation guarantee for the sampled value function}}
\setlength{\parindent}{0pt}
\setlength{\parskip}{5pt}
\raggedbottom
\newcommand{\pos}[1]{\bigl(#1\bigr)_+}
\newcommand{\E}{\mathbb E}
\newcommand{\Prob}{\mathbb P}
\newcommand{\R}{\mathbb R}
\newcommand{\lead}[1]{\par\smallskip\noindent\textbf{#1}\enspace}

\begin{document}
{\LARGE\bfseries Theorem 1. An approximation guarantee\\[2pt] for the sampled value function\par}
\medskip
A value function fitted from sampled states competes with the best uniform
approximation in its feature class. The guarantee holds simultaneously for
every optimizer, including when the fit is nonunique.

\section*{The sampled problem}
Let \(\mathcal S,\mathcal A\) be finite, nonempty state and action sets,
and let \(0<\alpha<1\) be the discount. Action \(a\) at state \(x\) costs
\(g_{x,a}\in\R\) and moves to \(y\) with probability \(p(x,y,a)\).
The Bellman operator for cost minimization is
\[
 (TJ)(x)=\min_{a\in\mathcal A}
   \left\{g_{x,a}+\alpha\sum_{y\in\mathcal S}p(x,y,a)J(y)\right\}.
\]
Write \(J^*\) for the optimal discounted cost-to-go and fix an optimal policy
\(\mu^*\). Let \(\nu\) be an initial probability distribution. For any
value vector \(f\), write \(\nu^\top f=\sum_x\nu_x f(x)\),
\(\|f\|_{1,\nu}=\sum_x\nu_x|f(x)|\), and \(\|f\|_\infty=\max_x|f(x)|\).

A feature map \(\Phi\) takes states into a real Hilbert space \(\mathcal H\).
For radii \(C,B\ge0\), the candidate functions are
\(\widetilde J_{\mathbf z,b}(x)=\langle\Phi(x),\mathbf z\rangle+b\),
with \(\|\mathbf z\|_{\mathcal H}\le C\) and \(|b|\le B\).
Draw \(X_1,\ldots,X_N\) independently from \(\pi_{\mu^*,\nu}\), the
discounted state distribution defined in \eqref{eq:occ}. Fit the parameters
by solving
\[
\begin{aligned}
 \operatorname{maximize}_{\mathbf z,b,s}\quad&
 \nu^\top\widetilde J_{\mathbf z,b}
       -\frac{2}{(1-\alpha)N}\sum_{i=1}^N s_i\\
 \text{subject to}\quad&
 \widetilde J_{\mathbf z,b}(X_i)\le (T\widetilde J_{\mathbf z,b})(X_i)+s_i,
 \quad s_i\ge0 \quad (1\le i\le N),\\
 &\|\mathbf z\|_{\mathcal H}\le C,\quad |b|\le B .
\end{aligned}
\tag{18}
\]
\textbf{Theorem 1 (strengthened form).} Let \(\epsilon>0\), \(0<\delta<1\) and \(N\ge1\).
Write \(K=\max_x\|\Phi(x)\|_{\mathcal H}\), \(\|g\|_\infty=\max_{x,a}|g_{x,a}|\), and
\[
 D=4CK(1+\alpha)+4B(1-\alpha)+2\|g\|_\infty,\qquad
 \Xi=D\bigl(1+\sqrt{2\ln(1/\delta)}\bigr).
\]
If \(N\ge\Xi^2/\epsilon^2\), then with probability at least \(1-\delta\),
every fitted value function \(\widehat J\) obtained by solving (18) satisfies
\begin{equation}
 \|J^*-\widehat J\|_{1,\nu}
 \le \frac{3+\alpha}{1-\alpha}
       \inf_{\|\mathbf z\|_{\mathcal H}\le C,\,|b|\le B}
       \|J^*-\widetilde J_{\mathbf z,b}\|_\infty
       +\frac{2\epsilon}{1-\alpha}.
 \tag{19$'$}\label{eq:main}
\end{equation}
A solution of (18) exists for every draw.

\newpage
\section*{Discounted occupancy and the Bellman residual}
Fix any policy \(\mu\), and set \(g_\mu(x)=g_{x,\mu(x)}\) and
\(P_\mu(x,y)=p(x,y,\mu(x))\). Its expected cost vector and its one-step
update of a candidate \(J\) are
\[
 J_\mu=\sum_{t=0}^\infty\alpha^tP_\mu^t g_\mu,\qquad
 T_\mu J=g_\mu+\alpha P_\mu J ,
\]
where \(P_\mu^0=I\), the identity matrix. The residual \(J-T_\mu J\)
measures how far the proposed cost exceeds its one-step prediction.

\textbf{Lemma (discounted occupancy).} For every \(\mu\) and \(\nu\) as above, the row vector
\begin{align}
 \pi_{\mu,\nu}^{\top}
 =(1-\alpha)\sum_{t=0}^{\infty}\alpha^t\nu^\top P_\mu^t
 =(1-\alpha)\nu^\top(I-\alpha P_\mu)^{-1}
 \tag{17}\label{eq:occ}
\end{align}
is a probability distribution. It is the unique real row vector satisfying
\begin{align}
 \pi_{\mu,\nu}^{\top}=(1-\alpha)\nu^\top+\alpha\pi_{\mu,\nu}^{\top}P_\mu .
 \tag{B}\label{eq:balance}
\end{align}
For every \(J\in\R^{\mathcal S}\), the expected initial value error satisfies
\begin{align}
 \nu^\top(J-J_\mu)
   =\frac{1}{1-\alpha}\,\pi_{\mu,\nu}^{\top}(J-T_\mu J).
 \tag{R}\label{eq:residual}
\end{align}

\lead{Proof. The distribution of discounted visits.}
The row vector \(\nu^\top P_\mu^t\) is the state distribution at time \(t\).
Indeed, multiplying a probability row by \(P_\mu\) preserves
nonnegativity and total mass. The coefficients \((1-\alpha)\alpha^t\)
are nonnegative and sum to one. Thus their mixture in \eqref{eq:occ}
converges coordinatewise and has total mass one. A state with zero initial
mass can have positive occupancy if the policy can reach it later.

\lead{The inverse and the balance equation.}
Every entry of \(P_\mu^t\) lies between zero and one, since each row is
a probability distribution. The matrix series \(\sum_{t\ge0}\alpha^tP_\mu^t\)
therefore converges absolutely in each entry. For the partial sum
\(D_n=\sum_{t=0}^{n}\alpha^tP_\mu^t\), multiplication on either side gives
\[
  (I-\alpha P_\mu)D_n
      =D_n(I-\alpha P_\mu)=I-\alpha^{n+1}P_\mu^{n+1}.
\]
The last term on the right tends to zero entrywise. There are only finitely
many states, so multiplication by \(I-\alpha P_\mu\) commutes with the
limit. The sum is consequently its two-sided inverse. This proves the
second formula in \eqref{eq:occ}.

\newpage
\section*{From sampled slacks to value error}
Multiplying (17) on the right by \(I-\alpha P_\mu\) gives the balance
equation \eqref{eq:balance}. Conversely, multiplying that equation by the
inverse uniquely determines \(\pi_{\mu,\nu}\). Since
\(J_\mu=g_\mu+\alpha P_\mu J_\mu\), subtraction gives
\[
 J-J_\mu=(I-\alpha P_\mu)^{-1}(J-T_\mu J).
\]
Averaging under \(\nu\) proves the signed identity \eqref{eq:residual}.
The inverse is entrywise nonnegative by its series on the previous page.
Writing \(u_+=\max\{u,0\}\) coordinatewise, we therefore also have
\[
 (J-J_\mu)_+\ \le\ (I-\alpha P_\mu)^{-1}(J-T_\mu J)_+ .
\]

\lead{Which residual does the fit pay for?}
Write \(r_\mu(J)=(J-T_\mu J)_+\) and \(r(J)=(J-TJ)_+\).
Since \(TJ\le T_\mu J\), we have \(r_\mu(J)\le r(J)\). In (18), the least
feasible slack at draw \(i\) is exactly \(r(J)(X_i)\). The objective strictly
decreases with each slack, so using these least slacks gives the reduced
objective \(\nu^\top J-2P_N r(J)/(1-\alpha)\), where
\[
 P_N f=\frac1N\sum_{i=1}^N f(X_i),\qquad Pf=\pi_{\mu^*,\nu}^\top f .
\]
An optimizer exists. Projecting \(\mathbf z\) orthogonally onto the span of
the finitely many \(\Phi(x)\) preserves every fitted value and weakly decreases
the norm. The closed ball in this span and the interval \([-B,B]\) are
compact. Both the reduced objective and the uniform approximation error
are continuous there, so their maximum and minimum, respectively, are attained.

\lead{A single deviation suffices.}
Suppose that \(P r_{\mu^*}(J)\le P_Nr_{\mu^*}(J)+\epsilon\) for every
admissible \(J\). Using
\(|u|=-u+2u_+\) with \(u=\widehat J-J^*\), the positive-part inequality,
and (17) with \(J_{\mu^*}=J^*\), gives
\begin{align*}
 \|J^*-\widehat J\|_{1,\nu}
 &\le \nu^\top(J^*-\widehat J)+\frac{2}{1-\alpha}P r_{\mu^*}(\widehat J)\\
 &\le \nu^\top(J^*-\widehat J)
        +\frac{2}{1-\alpha}\bigl(P_Nr(\widehat J)+\epsilon\bigr)\\
 &\le \nu^\top(J^*-V)+\frac{2}{1-\alpha}\bigl(P_Nr(V)+\epsilon\bigr)
\end{align*}
for every admissible \(V\). The last line uses optimality in the reduced
objective. It holds for each optimizer on the very same sample.

\newpage
\section*{A pointwise bound for the comparator}
Let \(V\) be any admissible comparison function. Each transition row is a
probability vector. Changing a value function by at most \(d\) everywhere
changes every one-step action value by at most \(\alpha d\). Taking the
minimum over actions preserves this bound. Hence
\(\|TV-TJ^*\|_\infty\le\alpha\|V-J^*\|_\infty\), and \(TJ^*=J^*\) gives
\[
 r(V)(x)\le |V(x)-J^*(x)|+|TJ^*(x)-TV(x)|
           \le (1+\alpha)\|J^*-V\|_\infty .
\]
This inequality holds at every sampled state, including repeated draws.
Also \(\nu^\top(J^*-V)\le\|J^*-V\|_\infty\). Substituting both bounds
in the last line on the previous page leaves the coefficient
\[
 1+\frac{2(1+\alpha)}{1-\alpha}=\frac{3+\alpha}{1-\alpha}.
\]
Taking the infimum over \(V\) yields \eqref{eq:main} on the stated
deviation event. It remains to bound that event's probability.

\section*{The size of the residual class}
Fix a policy \(\mu\). For \(\widetilde J=\widetilde J_{\mathbf z,b}\) put
\[
 d_\mu(x)=\Phi(x)-\alpha\sum_y P_\mu(x,y)\Phi(y),\qquad
 \|d_\mu(x)\|_{\mathcal H}\le K(1+\alpha).
\]
The norm bound follows because the transition weights sum to one. Then
\[
 r_\mu(\widetilde J)(x)=
 \pos{\langle d_\mu(x),\mathbf z\rangle+(1-\alpha)b-g_\mu(x)} .
\]
In particular \(0\le r_\mu(\widetilde J)\le CK(1+\alpha)+B(1-\alpha)
+\|g\|_\infty\le D/2\). The same envelope holds for \(r(\widetilde J)\)
by choosing a minimizing action at each state.

For a fixed sample \(X\), define the one-sided empirical Rademacher complexity
of a bounded function class \(\mathcal F\) by
\[
 \widehat R_N(\mathcal F;X)
 =\E_\sigma\sup_{f\in\mathcal F}\frac1N\sum_{i=1}^N\sigma_i f(X_i),
\]
where the independent signs \(\sigma_i\) take values \(+1,-1\) with equal
probability. For the linear term, Cauchy--Schwarz bounds the supremum by
\(C\|\sum_i\sigma_i d_\mu(X_i)\|/N\). The expected squared norm is
\(\sum_i\|d_\mu(X_i)\|^2\), since distinct signs have zero product in
expectation. Thus the complexity is at most \(CK(1+\alpha)/\sqrt N\).
The intercept similarly contributes at most \(B(1-\alpha)/\sqrt N\),
using \(\E|\sum_i\sigma_i|\le\sqrt N\).

\newpage
\section*{One event for all fitted values}
Let \(\mathcal F=\{r_{\mu^*}(\widetilde J_{\mathbf z,b}) :
\|\mathbf z\|\le C,\ |b|\le B\}\), and let \(-\mathcal F=\{-f:f\in\mathcal F\}\).
Taking the positive part after subtracting a fixed cost is 1-Lipschitz,
as is its negative. Applying either map at each state cannot increase
the one-sided complexity. To see this, condition on all signs except
one. If \(a_j\) is the sum of the other terms for candidate \(j\) and
\(u_j\) its remaining score, twice the conditional expected supremum is
\[
 \sup_{j,k}\{a_j+a_k+\varphi(u_j)-\varphi(u_k)\}
 \le \sup_{j,k}\{a_j+a_k+|u_j-u_k|\}.
\]
Interchanging \(j,k\) shows that the right side equals the expression for
the identity map. Repeating this for each coordinate removes the
transformations. The supremum of a sum is bounded by the sum of the
suprema, so the preceding page gives, for every \(X\),
\[
 \widehat R_N(-\mathcal F;X)
 \le \frac{CK(1+\alpha)+B(1-\alpha)}{\sqrt N}
 \le \frac{D}{2\sqrt N}.
\]

Set \(Z(X)=\sup_{f\in\mathcal F}(Pf-P_N f)\). It is measurable because
there are finitely many possible samples. For an independent copy \(X'\)
with empirical mean \(P_N'\), the supremum of expectations is bounded by the expectation of the supremum. Thus \(\E Z\le \E\sup_f(P_N'f-P_N f)\). Independently swapping \(X_i\)
and \(X_i'\) preserves their joint law and inserts independent signs into
the difference. Splitting the signed sum at the supremum and using
symmetry of the signs gives
\[
 \E Z\le 2\,\E_X \widehat R_N(-\mathcal F;X)\le D/\sqrt N .
\]

Changing one observation changes \(Z\) by at most \(2D/N\), since
\(|f|\le D\). The bounded-differences inequality for independent inputs
with coordinate changes at most \(c_i\) gives
\(\Prob(Z-\E Z\ge t)\le\exp(-2t^2/\sum_i c_i^2)\) for \(t>0\).
For completeness, expose the observations one at a time. The successive
increments of the conditional expectation of \(Z\) have conditional mean
zero and conditional range of width at most \(c_i\). For a variable in an
interval of width \(c_i\), the second derivative of its log moment-generating
function is a variance under exponential tilting, at most \(c_i^2/4\).
Integrating twice bounds each centered log moment-generating function
by \(\lambda^2 c_i^2/8\). Iterated conditioning, followed by Markov's
inequality with \(\lambda=4t/\sum_i c_i^2\), gives the stated tail.
For \(D>0\), take \(c_i=2D/N\). With probability at least \(1-\delta\),
\[
 Z\le\frac{D}{\sqrt N}+D\sqrt{\frac{2\ln(1/\delta)}{N}}
   =\frac{\Xi}{\sqrt N}\le\epsilon .
\]
For \(D=0\) all residuals vanish and the event has probability one.
This proves \eqref{eq:main}. \hfill\(\square\)

\newpage
\section*{A two-state example of discounted occupancy}
The states are \(s,t\), and the initial state is \(s\). From either state,
action \(A\) costs \(1\) and moves to \(s\) or \(t\) with equal probability.
Action \(B\) costs \(2\) and moves to \(t\) with certainty. The policy
takes \(A\) at \(s\) and \(B\) at \(t\). In state order \(s,t\),
\[
 P_\mu=\begin{pmatrix}1/2&1/2\\0&1\end{pmatrix},\qquad
 g_\mu=\begin{pmatrix}1\\2\end{pmatrix},\qquad
 \nu^\top=(1,0),\qquad \alpha=\tfrac12 .
\]
Writing \(u=\pi_{\mu,\nu}(s)\), the first coordinate of
\eqref{eq:balance} is \(u=1/2+u/4\). Thus the occupancy is
\((2/3,1/3)\). State \(t\) accounts for a third of discounted time even
though the initial distribution gives it zero mass.

At \(t\), the policy pays \(2\) forever, giving \(J_\mu(t)=4\). At \(s\),
the Bellman equation gives
\(J_\mu(s)=1+\tfrac14 J_\mu(s)+\tfrac14 J_\mu(t)\), so \(J_\mu(s)=8/3\).
Take the candidate \(J(s)=3,\ J(t)=2\). Its one-step predictions and
errors are
\[
 \begin{array}{c@{\qquad}r@{\quad}r@{\quad}r@{\quad}r}
   \text{State} & J_\mu & J & T_\mu J & J-T_\mu J\\ \hline
   s & 8/3 & 3 & 9/4 & \phantom{-}3/4\\
   t & 4   & 2 & 3   & -1
 \end{array}
\]
The residuals at the two states have opposite signs. Their contributions
partly cancel when averaged over discounted visits,
\[
 \pi_{\mu,\nu}^\top(J-T_\mu J)=\frac23\cdot\frac34+\frac13\cdot(-1)=\frac16 .
\]
Dividing by \(1-\alpha=1/2\) gives an initial error of \(1/3\). Direct
evaluation at the starting state gives the same answer,
\(J(s)-J_\mu(s)=3-8/3=1/3\).

\section*{Correction and comparison}
Appendix B.2 assigns a comparator \(V\) the slack \((TV-V)_+\).
For a single state with a self-transition, zero cost, \(\alpha=1/2\),
and \(V=1\), we have \(TV=1/2\). This candidate is admissible with
\(\Phi=0\), \(C=0\) and \(B=1\). The proposed slack is zero and violates
\(1\le1/2+s\). The least feasible slack is \((V-TV)_+=1/2\).

At the same sample size, \eqref{eq:main} implies the original guarantee. With probability at least \(1-\delta-\delta^4\), every optimizer satisfies
\begin{equation}
 \|J^*-\widehat J\|_{1,\nu}
 \le \frac{3+\alpha}{1-\alpha}
       \inf_{\|\mathbf z\|_{\mathcal H}\le C,\,|b|\le B}
       \|J^*-\widetilde J_{\mathbf z,b}\|_\infty+\frac{4\epsilon}{1-\alpha}.
 \tag{19}
\end{equation} The proof here uses the pointwise comparator
bound for the ordinary sup norm. The weighted comparator in Theorem~2
is a separate claim. We state the informative range \(0<\delta<1\);
for \(\delta\ge1\) the printed probability lower bound is nonpositive.

{\small
\textbf{Sources.}
\href{https://moallemi.com/ciamac/papers/kernel-adp-2012.pdf\#page=13}{Bhat, Farias, Moallemi and Zheng, Theorem~1 and Appendix~B}.
\href{https://github.com/nikhilpb/kadp/blob/main/lean/Theorem1.lean}{Lean theorem and optimizer existence},
\href{https://github.com/nikhilpb/kadp/blob/main/lean/Approximation.lean}{deterministic bound},
\href{https://github.com/nikhilpb/kadp/blob/main/lean/UniformDeviation.lean}{concentration},
\href{https://github.com/nikhilpb/kadp/blob/main/lean/Examples/Theorem1.lean}{slack example},
and \href{https://github.com/nikhilpb/kadp/blob/main/lean/Examples/Occupancy.lean}{occupancy example}.}
\end{document}
